\section{Introducción}

La gestión de memoria es una de las funciones fundamentales de un sistema operativo, ya que permite administrar eficientemente la memoria RAM para ejecutar múltiples procesos de manera concurrente.

En este informe se analizan distintos algoritmos de asignación de memoria implementados en Python mediante la clase \texttt{GestorMemoria}. Se estudian los siguientes métodos:

\begin{itemize}
    \item Particiones fijas
    \item Particiones dinámicas
    \item Particiones dinámicas con compactación
    \item Algoritmo de gemelos (Buddy System)
\end{itemize}

Además, se analizan los algoritmos de selección de bloques:

\begin{itemize}
    \item First Fit
    \item Best Fit
    \item Worst Fit
\end{itemize}

Finalmente, se explica detalladamente cómo cada algoritmo fue implementado dentro del código proporcionado.

%------------------------------------------------

\section{Objetivos}

\subsection{Objetivo General}

Analizar e implementar distintos algoritmos de gestión de memoria utilizados en sistemas operativos.

\subsection{Objetivos Específicos}

\begin{itemize}
    \item Comprender el funcionamiento de las particiones fijas y dinámicas.
    
    \item Analizar el problema de fragmentación interna y externa.
    
    \item Estudiar el funcionamiento de la compactación de memoria.
    
    \item Implementar y analizar el algoritmo de gemelos.
    
    \item Comparar los algoritmos First Fit, Best Fit y Worst Fit.
\end{itemize}

%------------------------------------------------

\section{Marco Teórico}

\subsection{Gestión de Memoria}

La gestión de memoria es el proceso mediante el cual el sistema operativo administra la memoria RAM disponible para los procesos.

Cada proceso necesita una cantidad determinada de memoria para ejecutarse, por lo que el sistema debe decidir:

\begin{itemize}
    \item Dónde almacenar cada proceso.
    \item Cómo reutilizar memoria liberada.
    \item Cómo minimizar desperdicio de memoria.
\end{itemize}

%------------------------------------------------

\subsection{Fragmentación}

La fragmentación representa memoria desperdiciada durante la asignación.

\subsubsection{Fragmentación Interna}

Ocurre cuando un proceso utiliza menos memoria que el tamaño del bloque asignado.

\textbf{Ejemplo:}

\begin{itemize}
    \item Bloque asignado: 256 KB
    \item Proceso: 200 KB
    \item Fragmentación interna: 56 KB
\end{itemize}

\subsubsection{Fragmentación Externa}

Ocurre cuando existen pequeños espacios libres dispersos en memoria que no pueden utilizarse eficientemente.

%------------------------------------------------

\section{Particiones Fijas}

\subsection{Definición}

En este esquema, la memoria se divide previamente en bloques de tamaño fijo.

Cada partición puede almacenar únicamente un proceso.

\subsection{Ventajas}

\begin{itemize}
    \item Implementación sencilla.
    \item Administración rápida.
\end{itemize}

\subsection{Desventajas}

\begin{itemize}
    \item Alta fragmentación interna.
    \item Limitación en el número de procesos simultáneos.
\end{itemize}

\subsection{Implementación}

Las particiones se crean en el constructor:

\begin{lstlisting}[language=Python]
if tipo_memoria == "Fijas":

    inicio = 0

    for tam in particiones_fijas:

        self.bloques.append(
            BloqueMemoria(inicio, tam)
        )

        inicio += tam
\end{lstlisting}

\subsection{Asignación de Procesos}

\begin{lstlisting}[language=Python]
bloque.libre = False

bloque.proceso = proceso

bloque.fragmentacion = (
    bloque.tamano - proceso.memoria
)
\end{lstlisting}

La fragmentación interna se calcula mediante:

\[
Fragmentacion = TamanoBloque - TamanoProceso
\]

%------------------------------------------------

\section{Particiones Dinámicas}

\subsection{Definición}

La memoria no se divide previamente. Los bloques son creados dinámicamente según las necesidades de cada proceso.

\subsection{Ventajas}

\begin{itemize}
    \item Menor fragmentación interna.
    \item Mejor aprovechamiento de memoria.
\end{itemize}

\subsection{Desventajas}

\begin{itemize}
    \item Produce fragmentación externa.
\end{itemize}

\subsection{Inicialización}

Inicialmente existe un único bloque libre:

\begin{lstlisting}[language=Python]
self.bloques.append(
    BloqueMemoria(0, total_ram)
)
\end{lstlisting}

\subsection{División de Bloques}

\begin{lstlisting}[language=Python]
nuevo = BloqueMemoria(
    bloque.inicio,
    proceso.memoria,
    libre=False,
    proceso=proceso
)
\end{lstlisting}

Cálculo del espacio restante:

\begin{lstlisting}[language=Python]
restante = (
    bloque.tamano - proceso.memoria
)
\end{lstlisting}

Creación del bloque libre restante:

\begin{lstlisting}[language=Python]
if restante > 0:

    libre = BloqueMemoria(
        nuevo.inicio + proceso.memoria,
        restante
    )

    self.bloques.insert(
        indice + 1,
        libre
    )
\end{lstlisting}

%------------------------------------------------

\section{Particiones Dinámicas con Compactación}

\subsection{Definición}

La compactación reorganiza los bloques ocupados para unir toda la memoria libre en una sola región continua.

\subsection{Ventajas}

\begin{itemize}
    \item Reduce fragmentación externa.
    \item Permite reutilizar memoria dispersa.
\end{itemize}

\subsection{Desventajas}

\begin{itemize}
    \item Alto costo computacional.
    \item Reubicación constante de procesos.
\end{itemize}

\subsection{Implementación}

Si no existe memoria suficiente:

\begin{lstlisting}[language=Python]
if not libres:

    if self.tipo_memoria == "Dinamicas Compactacion":

        self.compactar()
\end{lstlisting}

Compactación de bloques:

\begin{lstlisting}[language=Python]
for b in ocupados:

    nuevo = BloqueMemoria(
        inicio,
        b.tamano,
        libre=False,
        proceso=b.proceso
    )

    nuevos.append(nuevo)

    inicio += b.tamano
\end{lstlisting}

Creación del bloque libre final:

\begin{lstlisting}[language=Python]
libre_total = self.total_ram - inicio
\end{lstlisting}

%------------------------------------------------

\section{Algoritmo de Gemelos (Buddy System)}

\subsection{Definición}

El algoritmo de gemelos divide la memoria en bloques cuyo tamaño corresponde a potencias de 2.

\subsection{Funcionamiento}

\begin{enumerate}
    \item Se busca un bloque libre adecuado.
    
    \item Si el bloque es demasiado grande, se divide en dos mitades iguales.
    
    \item La división continúa hasta obtener el tamaño mínimo suficiente.
    
    \item Cuando un bloque se libera, puede fusionarse con su bloque gemelo.
\end{enumerate}

%------------------------------------------------

\subsection{Obtención de Potencia de 2}

\begin{lstlisting}[language=Python]
tam_requerido = 1

while tam_requerido < proceso.memoria:
    tam_requerido *= 2
\end{lstlisting}

%------------------------------------------------

\subsection{División de Bloques}

\begin{lstlisting}[language=Python]
while (
    bloque.tamano // 2 >= tam_requerido
):
\end{lstlisting}

Creación de bloques gemelos:

\begin{lstlisting}[language=Python]
bloque1 = BloqueMemoria(
    bloque.inicio,
    mitad
)

bloque2 = BloqueMemoria(
    bloque.inicio + mitad,
    mitad
)
\end{lstlisting}

%------------------------------------------------

\subsection{Fusión de Gemelos}

Verificación de buddies:

\begin{lstlisting}[language=Python]
if a.inicio ^ a.tamano == b.inicio:
\end{lstlisting}

Fusión:

\begin{lstlisting}[language=Python]
nuevo = BloqueMemoria(
    min(a.inicio, b.inicio),
    a.tamano * 2
)
\end{lstlisting}

%------------------------------------------------

\section{Algoritmos de Asignación}

\subsection{First Fit}

Selecciona el primer bloque disponible que sea suficientemente grande.

\subsubsection{Implementación}

\begin{lstlisting}[language=Python]
if self.algoritmo == "First Fit":

    bloque = libres[0]
\end{lstlisting}

\subsubsection{Ventajas}

\begin{itemize}
    \item Muy rápido.
    \item Bajo costo computacional.
\end{itemize}

\subsubsection{Desventajas}

\begin{itemize}
    \item Incrementa fragmentación externa.
\end{itemize}

%------------------------------------------------

\subsection{Best Fit}

Selecciona el bloque más pequeño posible que pueda almacenar el proceso.

\subsubsection{Implementación}

\begin{lstlisting}[language=Python]
bloque = min(
    libres,
    key=lambda b: b.tamano
)
\end{lstlisting}

\subsubsection{Ventajas}

\begin{itemize}
    \item Reduce desperdicio inmediato.
\end{itemize}

\subsubsection{Desventajas}

\begin{itemize}
    \item Genera bloques libres pequeños.
\end{itemize}

%------------------------------------------------

\subsection{Worst Fit}

Selecciona el bloque libre más grande disponible.

\subsubsection{Implementación}

\begin{lstlisting}[language=Python]
bloque = max(
    libres,
    key=lambda b: b.tamano
)
\end{lstlisting}

\subsubsection{Ventajas}

\begin{itemize}
    \item Mantiene bloques libres grandes.
\end{itemize}

\subsubsection{Desventajas}

\begin{itemize}
    \item Puede desperdiciar memoria.
\end{itemize}

%------------------------------------------------

\section{Análisis General de la Implementación}

\subsection{Diseño Modular}

La implementación utiliza programación orientada a objetos mediante:

\begin{itemize}
    \item Clase \texttt{GestorMemoria}
    \item Clase \texttt{BloqueMemoria}
\end{itemize}

Esto permite reutilizar y extender funcionalidades fácilmente.

%------------------------------------------------

\subsection{Lista de Bloques}

Toda la memoria es administrada mediante:

\begin{lstlisting}[language=Python]
self.bloques
\end{lstlisting}

Cada elemento representa una región continua de memoria.

%------------------------------------------------

\subsection{Liberación y Unión de Memoria}

Cuando un proceso termina:

\begin{lstlisting}[language=Python]
b.libre = True

b.proceso = None

b.fragmentacion = 0
\end{lstlisting}

Posteriormente se fusionan bloques contiguos:

\begin{lstlisting}[language=Python]
self.unir_libres()
\end{lstlisting}

%------------------------------------------------

\subsection{Ventajas de la Implementación}

\begin{itemize}
    \item Código modular y extensible.
    
    \item Fácil visualización del estado de memoria.
    
    \item Permite comparar múltiples algoritmos.
    
    \item Simula correctamente fragmentación y compactación.
\end{itemize}

%------------------------------------------------

\subsection{Limitaciones}

\begin{itemize}
    \item No implementa memoria virtual.
    
    \item No existe paginación.
    
    \item No implementa swapping.
    
    \item Compactación posee alto costo computacional.
\end{itemize}

%------------------------------------------------

\section{Conclusiones}

\begin{itemize}
    \item Las particiones fijas son simples, pero generan alta fragmentación interna.
    
    \item Las particiones dinámicas aprovechan mejor la memoria, aunque generan fragmentación externa.
    
    \item La compactación mejora el aprovechamiento de memoria al unir espacios dispersos.
    
    \item El algoritmo de gemelos proporciona divisiones y fusiones eficientes utilizando bloques potencia de 2.
    
    \item First Fit prioriza velocidad, Best Fit minimiza desperdicio inmediato y Worst Fit intenta preservar bloques grandes.
    
    \item La implementación desarrollada permite comparar diferentes estrategias de asignación de memoria de manera modular y flexible.
\end{itemize}

%------------------------------------------------

\section{Código Completo Implementado}

\begin{lstlisting}[language=Python]
from models.bloque_memoria import BloqueMemoria


class GestorMemoria:

    def __init__(
        self,
        total_ram,
        tipo_memoria,
        algoritmo,
        particiones_fijas=None
    ):

        self.total_ram = total_ram

        self.tipo_memoria = tipo_memoria

        self.algoritmo = algoritmo

        self.bloques = []

        if tipo_memoria == "Fijas":

            inicio = 0

            for tam in particiones_fijas:

                self.bloques.append(
                    BloqueMemoria(inicio, tam)
                )

                inicio += tam

        else:

            self.bloques.append(
                BloqueMemoria(0, total_ram)
            )

    def asignar(self, proceso):

        libres = [
            b for b in self.bloques
            if b.libre and b.tamano >= proceso.memoria
        ]

        if not libres:

            if self.tipo_memoria == "Dinamicas Compactacion":

                self.compactar()

                libres = [
                    b for b in self.bloques
                    if b.libre and b.tamano >= proceso.memoria
                ]

            if not libres:
                return False

        if self.tipo_memoria == "Gemelos":

            bloque = min(
                libres,
                key=lambda b: b.tamano
            )

            tam_requerido = 1

            while tam_requerido < proceso.memoria:
                tam_requerido *= 2

            while (
                bloque.tamano // 2 >= tam_requerido
            ):

                mitad = bloque.tamano // 2

                indice = self.bloques.index(bloque)

                bloque1 = BloqueMemoria(
                    bloque.inicio,
                    mitad
                )

                bloque2 = BloqueMemoria(
                    bloque.inicio + mitad,
                    mitad
                )

                self.bloques.pop(indice)

                self.bloques.insert(indice, bloque2)

                self.bloques.insert(indice, bloque1)

                bloque = bloque1

            bloque.libre = False

            bloque.proceso = proceso

            bloque.fragmentacion = (
                bloque.tamano - proceso.memoria
            )

            return True

        if self.algoritmo == "First Fit":

            bloque = libres[0]

        elif self.algoritmo == "Best Fit":

            bloque = min(
                libres,
                key=lambda b: b.tamano
            )

        else:

            bloque = max(
                libres,
                key=lambda b: b.tamano
            )

        if self.tipo_memoria == "Fijas":

            bloque.libre = False

            bloque.proceso = proceso

            bloque.fragmentacion = (
                bloque.tamano - proceso.memoria
            )

            return True

        nuevo = BloqueMemoria(
            bloque.inicio,
            proceso.memoria,
            libre=False,
            proceso=proceso
        )

        restante = (
            bloque.tamano - proceso.memoria
        )

        indice = self.bloques.index(bloque)

        self.bloques.pop(indice)

        self.bloques.insert(indice, nuevo)

        if restante > 0:

            libre = BloqueMemoria(
                nuevo.inicio + proceso.memoria,
                restante
            )

            self.bloques.insert(
                indice + 1,
                libre
            )

        return True
\end{lstlisting}

\section{Repositorio del Proyecto}

El código fuente completo del proyecto se encuentra disponible en GitHub:

\begin{center}
\href{https://github.com/Elemauz/Gestion-de-Memoria}
{Repositorio Oficial del Proyecto}
\end{center}